We built a system that takes one ECG beat and returns its class, normal (N), supraventricular ectopic (S), ventricular ectopic (V) or fusion (F), a saliency map of the parts of the beat behind the prediction, and a short note for a clinician that is checked against the model output. With it we asked whether synthetic beats of the rare classes from a conditional VAE or a WGAN-GP help a model find them in patients it has never seen. Every choice was fixed on validation patients by rules written down in advance, and the 22 test patients of the MIT-BIH Arrhythmia Database were read once. Generative augmentation was detectably better than no augmentation in 0 of 4 pre-registered comparisons, and the best pre-registered model reached a rare-class macro-F1 (the mean of S and F F1) of 0.136. The cnn without augmentation, 0.088 on new patients, scored 0.777 to 0.837 with the same patients in training and test. Follow-up analyses, each registered before it ran, found what moves the score: each beat seen relative to the patient's own normal beat (0.088 to 0.371), prematurity judged against the recording's mean heart period, no S calls during atrial fibrillation, and ten runs averaged; together they reached 0.587 on the test records, close to the best published inter-patient result we could check (0.582). That score is optimistic after repeated reads of the test records: on 977 unused patients the pipeline's S F1 was 0.294, against 0.647 on the test records. Patient variation, not class rarity, limits these models; more training patients with rare beats, larger models and more runs are needed to establish how much augmentation can add.
